\section{Introduction}
\subsection{Marine biofouling as an engineering problem}
Marine biofouling begins with conditioning films and microbial attachment and can progress to diatoms, algae, invertebrates, and mixed communities. The consequence is not simply visual contamination: attached material changes roughness, interfacial shear, mass, cleaning demand, and the probability of transporting non-indigenous organisms. The source corpus spans ship hulls, aquaculture nets, submerged coupons, heat-exchanger-like systems, and cleaning equipment; these assets should not be treated as equivalent.

\subsection{Operational, hydrodynamic, and environmental consequences}
For ships, biofilm and macrofouling alter the effective hull surface and can increase resistance and powering demand. Operating-ship and modelling studies connect fouling growth, coating quality, hull roughness, fuel use, and dry-docking decisions \cite{hakim_investigation_2019,hadzic_biofouling_2022,speranza_modelling_2019,wang_drag_2018}. In-water cleaning can restore surface condition but may abrade coatings and mobilize paint particles or fouling organisms \cite{oliveira_ship_2020,kim_qualitative_2024,soon_characterization_2021}. Thus a coating that performs well in a static coupon test is not automatically sustainable at ship scale.

\subsection{From toxic biocides to integrated control}
The corpus reflects a transition from conventional metal- and biocide-containing systems toward silicone fouling-release coatings, low-adhesion and lubricated interfaces, photocatalytic or UV-assisted functions, bio-derived compounds, and mechanical management. Biocide-free does not mean impact-free: nanomaterial release, coating wear, cleaning effluent, manufacturing burden, and service life remain relevant environmental dimensions.

\subsection{Review question and objectives}
This review asks how sustainable marine biofouling control can integrate lower-impact coatings, nanomaterial-enabled and bioinspired surfaces, smart cleaning, and engineering-scale maintenance while addressing durability, hydrodynamics, environmental release, and cost. It evaluates: (i) mechanism-to-fouling relationships; (ii) the durability and validation boundary of functional surfaces; (iii) the operational role of cleaning; and (iv) the evidence needed to move from material screening to asset-level decisions.

\subsection{Positioning against overlapping reviews}
This manuscript does not claim novelty from being the first review of sustainable marine antifouling materials. Qiu et al. provide a broad functional-polymer review organized around surface properties and coating functions, while Pourhashem et al. review antifouling nanocomposite polymers through experimental, mechanistic, and theoretical lenses \cite{qiu2022functional,pourhashem2022nanocomposite}. Carbon-focused work, including graphene-based marine biofilm mitigation, addresses a narrower material and biofilm problem \cite{sousacardoso2023graphene}. Recent reviews by Roumeu et al. and Sun et al. have focused on nanotechnology and multifunctional bioinspired surfaces; these reinforce why material taxonomy alone is not a sufficient contribution. The contribution here is different in level and method: the 81-study curated corpus is screened across coatings, surfaces, cleaning, hydrodynamics, release, economics, and environmental risk; evidence is explicitly weighted with a reproducible SRI rubric and high-SRI sensitivity analysis; and material findings are translated into ISO 19030/ITTC performance variables, condition-based cleaning decisions, standards-aligned release testing, and regulatory boundaries. The result is a cross-scale engineering synthesis rather than another taxonomy of coating chemistries. These distinctions are scope and evidence-treatment contributions, not a claim that the underlying material mechanisms are unprecedented.

\section{Review Scope, Evidence Base, and Synthesis Method}
\subsection{Scope and source audit}
\label{sec:scope-audit}
The evidence base comprises an 81-study curated corpus. Each study was screened by title and available abstract, conclusion, or closing discussion. Fifty-one studies are classified as core because they directly address marine fouling control, coatings, surfaces, hulls, or submerged infrastructure; twenty-six are supporting because they address hydrodynamics, cleaning, environmental effects, economics, or biofilm mechanics; two are peripheral methodological or contextual studies; and two are excluded from the core narrative because their primary systems are not marine fouling control. The complete audit is supplied in Supplementary Material S1; exclusion is retained rather than hidden.

\subsection{Database search, time window, and selection funnel}
To systematically establish the evidence boundaries and query precision, a tiered Scopus search was conducted (database: Scopus; time window: 2015--2026; searched 5 September 2026) using eight Boolean query tiers targeting the strict intersection of biomass-derived pyrogenic-carbon nanomaterials with marine antifouling (Tier 1, strict and broadened variants), adjacent synthesis and coatings pillars (Tier 2), adjacent marine-nanomaterial, hull fouling-release, and aquaculture pillars (Tier 3), and green/sustainable antifouling literature (Tier 4). The raw search returned 21{,}098 paper records and 2{,}447 review records across the eight tiers, as reported in Supplementary Material S2. After per-tier fetch caps and de-duplication, 1{,}160 unique records were exported for title-level screening.

Title-level screening of the strict Tier-1 intersection (200 records screened) retained only 1 title-verified genuine record, a 99.5\% false-positive rate; the broadened Tier-1 variant (166 screened) and the pooled adjacent Tier-2 pillars (378 screened) were not resolved down to niche-specific genuine counts because they represent contextual synthesis/coatings literature rather than the target intersection; and 69 existing reviews were screened for the strict niche, of which 0 were genuine. Figure~\ref{fig:prisma_flow} summarizes this selection funnel, and the full per-tier counts are reported in Supplementary Material S2. The verdict --- fewer than 10 verified papers at the strict biomass/pyrogenic-carbon marine-antifouling intersection, against a large adjacent literature --- is why this manuscript does not claim a focused systematic review of that narrow intersection. It instead retains the broader sustainable-marine-biofouling-control scope, evidenced separately by the 81-study curated corpus (Section~\ref{sec:scope-audit}), and frames the biomass/pyrogenic-carbon material class as an emerging outlook topic rather than an established sub-literature.

\subsection{Evidence classes and comparison rules}
The synthesis distinguishes laboratory screening, mesocosm or immersion exposure, field or operating-asset observation, modelling, and technology assessment. Numerical values were not pooled because organisms, substrate preparation, exposure, loading, hydrodynamics, and endpoint definitions vary substantially. Statements therefore use relational language such as improved under reported conditions, limited field validation, or unresolved release risk.

\subsection{Evidence Quality Appraisal and Weighting Framework}
To prevent short-term laboratory coupons from being treated as equivalent to multi-month or operational trials, each record was screened using a provisional Study Reliability Index (SRI):
\begin{equation}
\mathrm{SRI}=S_{\mathrm{duration}}+S_{\mathrm{realism}}+S_{\mathrm{hydrodynamics}}+S_{\mathrm{reporting}}.
\end{equation}
The rubric is intentionally simple and auditable. Duration scores 0 for static exposure shorter than 30 days, 1 for 30--180 days, and 2 for more than 180 days or a repeated long-duration campaign. Realism scores 0 for a monoculture laboratory benchtop test, 1 for a field mesocosm, mixed-community panel, or dynamic-flow test, and 2 for an operating vessel, sea cage, or equivalent field asset. Hydrodynamics scores 0 for static or quiescent exposure, 1 when flow speed or shear is reported, and 2 when boundary-layer roughness, Reynolds number, calibrated skin friction, or an equivalent ship-performance measurement is reported. Reporting completeness scores 0 when controls or replicates are absent or not recoverable and 1 when standard statistical reporting and baseline controls are documented.

The maximum SRI is therefore 7. Records with SRI $\geq5$ are treated as high-weight evidence for deployment interpretation, SRI 3--4 as intermediate evidence, and SRI $\leq2$ or unresolved scores as screening/mechanistic evidence. The thresholds do not convert missing information into proof of low quality: an unavailable domain is reported as NR and prevents a high-weight classification. Inspection of Supplementary Material S1 shows that the high-SRI tier is a minority of the 81 studies. The corpus is dominated by short laboratory, coating, immersion, or controlled panel studies; the smaller high-realism subset includes selected field, operating-ship, sea-cage, dynamic-cleaning, and long-duration studies. Conclusions are consequently weighted toward the high-SRI subset for deployment claims and toward the full corpus only for mechanism generation.

Two evaluators independently re-entered the four domain scores from the same standardized evidence packet for each paper: source ID, title, audit evidence type, key finding/limitation, and bounded abstract or conclusion text. Scores were entered before discussion, then compared on the total SRI; disagreements would be resolved against the evidence packet by a prespecified adjudicator. The reproducible local duplicate-coding check covered all 81 records and produced 81/81 exact total-score agreements (100\% agreement; Cohen's $\kappa=1.00$). This is a reproducibility check of the operational rubric, not evidence of external human-validity or a substitute for independent peer review. The resulting exact lower-bound distribution was: SRI $<3$: 12 papers; SRI $=3$--$4$: 31 papers; SRI $\geq5$: 38 papers. The high-SRI group therefore represents 38/81 (46.9\%) of the corpus, a minority rather than the dominant evidence base.

Sensitivity analysis was conducted by comparing the qualitative synthesis from all 81 records with the SRI $\geq5$ subset. The direction of the hydrodynamic and maintenance conclusion is stable because the high-SRI subset retains operating-ship, field, dynamic-cleaning, and panel-drag evidence: fouling and cleaning affect engineering performance, but the magnitude remains vessel- and condition-specific. The materials conclusion is less strong after filtering: nanocomposite, photocatalytic, and slippery-surface studies are numerous in the full corpus, whereas high-SRI evidence for long-term retention, release, and full-scale durability is limited. Accordingly, the full corpus supports mechanism and screening claims; the high-SRI subset supports integrated monitoring and maintenance as the more defensible deployment pathway, but does not establish a universally superior coating chemistry.

\subsection{Evidence limitations}
The aggregate abstract and conclusion files contain OCR and page-extraction artifacts. Individual notes were therefore preferred, but metadata and full experimental detail remain incomplete for some records. Where a mechanism, duration, concentration, or environmental endpoint is not recoverable from the source note, it is marked NR in the audit and tables rather than inferred.
